\documentclass[../main.tex]{subfiles}

\begin{document}

\section{Conclusions}\label{sec:conclusions}

This work transcribed the minimum-fuel powered-descent and pinpoint-landing optimal control problem into a finite-dimensional NLP via trapezoidal direct collocation, and solved it with \texttt{fmincon}/SQP. The main findings are as follows.

\textbf{Direct vs.\ indirect transcription.}\ Compared to the indirect-shooting approach of Homework~1, direct collocation requires no analytical derivation of costate equations or transversality conditions, accepts path constraints (glide-slope, thrust bounds) natively as inequality constraints, and is robust to a coarse initial guess. The cost is a higher-dimensional NLP ($350$ variables here) and the loss of the structured analytical insight that the bilinear-tangent law provided in the indirect formulation.

\textbf{Bang-bang thrust structure.}\ The fuel-optimal solution exhibits the classical bang-bang behaviour expected for problems with linear-in-control dynamics and a control-magnitude bound: $\|\mathbf{T}\|$ saturates at $T_{\max}$ over the deceleration phase and drops to a low value during the soft-landing approach. This is the direct-transcription analogue of the maximum-principle prediction.

\textbf{Sensitivity to flight time.}\ A $\pm 5\%$ sweep around the nominal $t_f = 38$~s revealed monotonic growth of fuel consumption with flight time (longer hover $\Rightarrow$ greater gravity losses), with a $\sim 1.2\%$ propellant variation across the window. The lower bound on $t_f$ is set implicitly by the thrust-magnitude constraint.

\textbf{Active constraints at the optimum.}\ Two inequalities are active over substantial portions of the solution: the upper thrust bound $\|\mathbf{T}\| = T_{\max}$ and the glide-slope cone in the late descent. By the KKT complementarity conditions the corresponding multipliers are positive, identifying these as the binding constraints that shape the trajectory.

\textbf{ZOH transcription (Task 2).}\ The optional ZOH discretisation of Section~\ref{sec:task2} was implemented in three complementary forms---a nonlinear ZOH with RK4 propagation, the literal Appendix~A construction (LTV-linearised + SCvx) solved by \texttt{fmincon}/SQP, and the same SCvx outer loop with the inner sub-problem cast as a Second-Order Cone Program and solved by \texttt{YALMIP}+\texttt{ECOS}. All three confirm the trapezoidal baseline at the engineering level: final masses agree within a few kilograms ($\sim 0.3\%$ of the propellant). The forward-integration test exposes the lower convergence order of the trapezoidal rule, with the ZOH--RK4 variant dropping the per-node fidelity error by five orders of magnitude and sitting close to the \texttt{ode45} tolerance floor. Warm-started from the trapezoidal solution and stabilised by an adaptive trust-region ratio test, both SCvx variants converge to the same local optimum. The conic variant exposes the SOCP structure of the inner step directly and delivers a per-step fidelity that is an order of magnitude tighter than the \texttt{fmincon}/SQP path, demonstrating the convex-sub-problem framework underlying production guidance algorithms (GFOLD, lossless convexification).

\textbf{Non-dimensionalisation.}\ Following the same approach used in Homework~1, the OCP is solved in non-dimensional variables (Section~\ref{sec:nondim}), with $L_{ref} = y_0$, $a_{ref} = g$ and the single residual parameter $V_c = V_{ref}/c \approx 0.078$. This shrinks the dynamic range of the decision vector from four orders of magnitude to one and is essential for the LTV linearisation to be well-conditioned: the Jacobians become uniformly $O(1)$ and \texttt{fmincon}/SQP makes appreciable progress per iteration.

\textbf{Limitations and roadmap.}\ The model neglects aerodynamic drag, vehicle attitude dynamics, and any throttling rate limits, and the flight time is fixed. Natural extensions remaining beyond the scope of this report are:
\begin{itemize}
    \item A \textbf{free-time variant} that lifts the fixed-$t_f$ assumption and allows the solver to find the absolute fuel-optimal flight time subject to the active thrust constraint, e.g.\ as in lossless-convexification formulations.
    \item \textbf{Lossless convexification} of a non-zero lower thrust bound $T_{\min} > 0$, which is non-convex in $(\mathbf{T})$: with the conic SCvx infrastructure already in place (variant~(c)), the next step is to introduce the slack-variable trick of Açikmeşe and Ploen to recover a convex SOCP in the new state, and verify that the original constraint is recovered at the optimum.
    \item Inclusion of \textbf{aerodynamic drag} and \textbf{attitude dynamics}, moving from a 3-DoF point mass to a 6-DoF rigid body with angle-of-attack and angle-of-sideslip dependence.
\end{itemize}

\end{document}
